\section{Conclusions}
\label{sec:conclusion}

AQSE connects versioned sensor-feature representations to an eight-qubit trainable fidelity kernel and classical prediction interfaces. A repository-preregistered experiment evaluated its harmonic-feature TQK--SVC component across 30 newly generated sensor corpora. The mean balanced accuracy was 0.9528 for TQK and 0.9569 for RBF-SVC. The paired difference was $-0.0042$ with a 95\% bootstrap interval of $[-0.0222,0.0125]$. The primary predictive-advantage criterion was not met.

QNG changed the recorded kernel spectra, but those changes do not establish a benefit in held-out prediction. Selected-kernel geometry was distinct from rank-one and identity limits, while the exploratory spectral--performance associations remained unresolved. The TRAIN-label control retained validation supervision and did not meet the chance-compatibility check for TQK and RBF. It therefore cannot serve as a complete label-null validation.

The contribution is a traceable framework and a bounded comparative evaluation, not evidence of quantum hardware advantage or field readiness. The study separates representational change from predictive benefit and preserves the negative and inconclusive findings. Future matched ablations and physical-sensor studies can test extensions without altering this frozen evidence.
